% 18 November 2026; prepared for Edward Baker with OpenAI Codex assistance.
% Actual model/effort unavailable. Source: CM Note27 and E25 assessment.
\section{Compact source acquisition and a local stopping theorem}
\label{sec:compact}

The positive-time source of the preceding section supplies both a joint
estimate and pre-existing exterior reality. A compact application has a
different logical requirement: it must acquire the needed information
without assuming that the target gap remains positive on the interval
being tested. A separating contour gives such a source interface because
its count may remain two through a target merger. We first derive that
interface for trial pair coordinates, then state the stopping theorem and
its full error requirements \cite{BakerCM27}.

\subsection{Complete moments at trial coordinates}

Fix a compact real time interval and a positively oriented circle
$\Gamma=\{|z-a|=r\}$, with real $a$. Suppose $H$ is nonzero on a
neighborhood of $\Gamma$ throughout the interval and counts exactly two
occurrences inside. These hypotheses allow a double target or a conjugate
nonreal pair. No assertion about exterior reality is made. The canonical
order-one product and its locally uniform logarithmic derivative are
used on this contour. Let $B=H_z/H$ and
\[
 A_j=\frac1{2\pi\ii}\oint_\Gamma(z-a)^jB(t,z)\dd z,
 \qquad c=a+A_1/2,\qquad D=A_2/2-(A_1/2)^2.
\]
These are the actual coordinates. For separate trial real coordinates
$(\widehat c,\widehat D)$ put
$Q_{\!\widehat c,\widehat D}(z)=(z-\widehat c)^2-\widehat D$.
Both roots of this trial quadratic are required to lie inside $\Gamma$;
they need not coincide with the actual targets.

\begin{proposition}[Trial-coordinate exterior traces]
\label{cmp:trial}
For every integer $p\ge1$,
\begin{equation}
 S_p^\Gamma(t;\widehat c,\widehat D)
 =-\frac1{2\pi\ii}\oint_\Gamma
       \frac{B(t,z)}{Q_{\!\widehat c,\widehat D}(z)^p}\dd z
 =\sum_{\rho\text{ outside }\Gamma}
       \bigl((\rho-\widehat c)^2-\widehat D\bigr)^{-p}.
 \label{cmp:trace}
\end{equation}
The sum includes the complete exterior with multiplicity. It converges
absolutely and locally uniformly on compact trial sets with strict
contour separation. At the actual $(c,D)$ it equals $S_p$.
\end{proposition}
\begin{proof}
Set $f=Q_{\!\widehat c,\widehat D}^{-p}$. All its poles are inside
$\Gamma$, and $f=O(z^{-2p})$ at infinity. For an actual interior root
$\rho$, the rational function $f(z)/(z-\rho)$ has all poles inside and
decays as $O(z^{-2p-1})$, so its contour integral is zero. Coincident or
multiple poles cause no difficulty. For an exterior root the residue
outside is $f(\rho)$, and the same decay gives the interior integral
$-f(\rho)$. The constant term and genus-one constant regularizers of the
logarithmic derivative integrate to zero since $\oint f=0$.
The complete product converges normally on the cleared contour. Its
exterior contributions to \eqref{cmp:trace} are dominated by a multiple
of $\sum_\rho(1+|\rho|)^{-2}$, finite for the order-one canonical divisor;
compact trial sets permit the same dominating tail. Thus termwise
contour passage gives the result. This argument extends the rational
trace of Proposition~\ref{prop:trace} without identifying trial and actual
pair coordinates.
\end{proof}

Define the complete signed trial forcing
\begin{equation}
 \begin{split}
 \mathcal F_\Gamma&=12S_2^\Gamma+16\widehat D S_3^\Gamma,\\
 K_\Gamma&=2\widehat D^2\mathcal F_\Gamma
 =-\frac1{2\pi\ii}\oint_\Gamma
 \left(\frac{24\widehat D^2}{Q^2}
       +\frac{32\widehat D^3}{Q^3}\right)B(t,z)\dd z .
 \end{split}
 \label{cmp:forcing}
\end{equation}
Here and below $Q$ in a trial integral denotes the trial quadratic.
Conjugation makes these quantities real for real trial coordinates.
Unlike a positive-head estimate, \eqref{cmp:forcing} does not discard
any possibly nonreal exterior terms.

At the actual coordinates the local pair law and the first hierarchy
row give, in backward time $u=t_*-t$,
\begin{equation}
 D_u=-2+4DM,\qquad M_u=-2M^2+\mathcal F,
 \qquad \mathcal F=12S_2+16DS_3.
 \label{cmp:local-system}
\end{equation}
The cancellation proof of Theorem~\ref{thm:hierarchy} uses finite rational
identities before complete contour passage; neither step uses positivity
of exterior summands. It therefore applies to the present strictly
separating local tube, with its own compact contour bounds in place of
the strong-collar constants. These hypotheses must hold on the interval
at issue. An earlier positive-time certificate does not supply them on
an enlarged negative-time interval.

\subsection{Stopping without an assumed positive trajectory}

Suppose unconditional contour information gives
$|c(t)-a|\le C$ and $D(t)\le D_{\max}$ throughout the desired interval.
Its lower enclosure for $D$ is allowed to cross zero. Choose
$0<d_{\mathrm{stop}}<d_0$ and assume an independent initial enclosure
$D(0)\ge d_0$, $V(0)\le v_0$, where $V=1-2DM$ uses backward time.

\begin{theorem}[A conditional source bound prevents a lower exit]
\label{cmp:stop}
Assume the count-two contour and its clearance hold for
$0\le u\le\Delta$. Suppose all trial quadratics with
$|\widehat c-a|\le C$ and
$d_{\mathrm{stop}}\le\widehat D\le D_{\max}$ lie inside the contour and
\begin{equation}
 K_\Gamma(t_* -u;\widehat c,\widehat D)\ge\kappa
 \label{cmp:kfloor}
\end{equation}
uniformly on that trial set. Put
\[
 v=\max\{0,v_0,\sqrt{\max(0,1-\kappa)}\}.
\]
If
\begin{equation}
 d_0-2v\Delta-d_{\mathrm{stop}}>0,
 \label{cmp:exit-margin}
\end{equation}
then the actual pair satisfies $D(u)>d_{\mathrm{stop}}$ throughout the
interval. Exterior zeros need not be real.
\end{theorem}
\begin{proof}
Until a first lower exit $D=d_{\mathrm{stop}}$, the actual coordinates
belong to the trial set. Thus $K\ge\kappa$ there, and
\eqref{cmp:local-system} gives
$D_u=-2V$ and $V_u=(1-V^2-K)/D$. For $w=V-v$,
\[
 w_u+\frac{V+v}{D}w\le\frac{1-v^2-\kappa}{D}\le0.
\]
On every compact pre-exit interval the integrating factor is positive,
so $V\le v$ and $D\ge d_0-2vu$. At a hypothetical first exit,
\eqref{cmp:exit-margin} gives $D>d_{\mathrm{stop}}$, a contradiction.
The other coordinate and contour exits are already excluded by the
unconditional assumptions. No division by $D$ at a merger is used.
\end{proof}

The distinction between the trial statement and an unconditional positive
forcing floor is substantive. On a regular counted tube $\mathcal F$
is finite at a target merger, whereas $K=2D^2\mathcal F$ vanishes there.
Hence an unconditional $K\ge k>0$ already excludes a merger by continuity
from a positive start. Such a floor can be a valid source certificate,
but is not a weaker substitute for the sought positivity. Theorem
\ref{cmp:stop} instead proves the necessary interval of positivity from
an estimate on a conditional trial region.

\subsection{A complete analytic error contract}

Let $G$ be any fixed nonzero scalar multiple of $H$. Suppose source
approximants $h,h_1$ have uniform errors $e_0,e_1$ for $G,G_z$ on the
whole time-contour cell, and $|h|\ge b_0>e_0$, $|h_1|\le b_1$. Then
\begin{equation}
 \left|B-\frac{h_1}{h}\right|
 \le e_B:=\frac{e_1}{b_0-e_0}
          +\frac{b_1e_0}{b_0(b_0-e_0)}.
 \label{cmp:quotient-error}
\end{equation}
Indeed, subtract the quotients and bound the numerator error over $G$
and the reciprocal-denominator error separately. If
$|\widehat c-a|\le C$, $|\widehat D|\le D_{\max}$,
$0\le C<r$, $D_{\max}\ge0$ and
$q=(r-C)^2-D_{\max}>0$, integration around the circle gives the source
parts of the moment and forcing errors:
\begin{align}
 E_M&\le re_B/q,&
 E_{A_j}&\le r^{j+1}e_B,\notag\\
 E_K&\le re_B\left(24D_{\max}^2/q^2+32D_{\max}^3/q^3\right).
 \label{cmp:kernel-errors}
\end{align}
All parameter dependence belongs to the trial enclosure; discarding a
proved time-state correlation by replacing it with independent boxes is
an additional approximation, not part of these inequalities.

An annular logarithm gives an alternative interface when
$\mathcal L=\Log(G/(z-a)^2)$ is single valued on a cleared annulus.
Since $B=\mathcal L_z+2/(z-a)$ and the rational second term has zero
integral against $Q^{-p}$, integration by parts gives
\begin{align}
 S_p^\Gamma&=-\frac{p}{2\pi\ii}\oint_\Gamma
                  \mathcal L\frac{Q'}{Q^{p+1}}\dd z,\notag\\
 K_\Gamma&=-\frac1{2\pi\ii}\oint_\Gamma\mathcal L Q'
              \left(\frac{48\widehat D^2}{Q^3}
                    +\frac{96\widehat D^3}{Q^4}\right)\dd z.
 \label{cmp:log-kernel}
\end{align}
A constant may be subtracted from $\mathcal L$ without changing either
integral. A uniform log error $e_{\mathcal L}$ gives
\begin{align}
 E_M&\le 2r(r+C)e_{\mathcal L}/q^2,\notag\\
 E_K&\le 2r(r+C)e_{\mathcal L}
        \left(48D_{\max}^2/q^3+96D_{\max}^3/q^4\right).
 \label{cmp:log-errors}
\end{align}
For matching log branches and $|G-h|\le e_0<b_0\le|h|$, the log increment
has modulus at most $-\log(1-e_0/b_0)$ by its convergent power series.
Branch matching and annular clearance remain explicit obligations.

Finite contour quadrature has its own error. If the parameterized
integrand is holomorphic for $|\Im\theta|\le s$ and bounded there by
$A_s$, its normalized $N$-node trapezoidal error is at most
\begin{equation}
 \frac{2A_s}{\e^{Ns}-1}.
 \label{cmp:trapezoid}
\end{equation}
To prove this, shift the Fourier coefficient contour to the appropriate
strip edge, obtaining $|a_k|\le A_s\e^{-|k|s}$. The trapezoidal mean
aliases only the coefficients $a_{jN}$; summing their geometric bound
gives \eqref{cmp:trapezoid}. The source and every trial denominator must
be cleared throughout that complex strip, not only at real nodes.

The budgets $e_0,e_1$ must include complete source evaluation, source
integration tails, lattice or quadrature errors, coefficient reconstruction,
arithmetic, spatial approximation and whole-cell time variation. Time
remainders may use $\partial_t^jG=(-1)^j\partial_z^{2j}G$ with actual
Cauchy bounds. Time samples alone do not supply uniform estimates, and
an asymptotic error cannot be differentiated without an analytic residual.
The complete contour representation requires no separately omitted zero
tail after the full source error has been paid.

A negative-time source formula is available analytically. For $s\ge0$,
\begin{equation}
 H(-s,z)=\pi^{-1/2}\int_\R\e^{-y^2}
                         H(0,z+2\sqrt{s}\,y)\dd y.
 \label{cmp:negative-gaussian}
\end{equation}
Insert the complete Fourier representation of $H(0,\cdot)$; its positive
source and fixed $\Im z$ dominate the double integral. The Gaussian
Fourier integral supplies the multiplier $\e^{-s u^2}$, proving the
formula by Fubini. The real spatial shift in \eqref{cmp:negative-gaussian}
differs from the imaginary shift at positive heat time. Consequently an
existing positive-time xi quadrature estimate cannot be extrapolated to
negative time without new source and tail bounds.

\subsection{What a constant barrier can and cannot retain}

The stopping theorem preserves coupling in its source kernel, but its
single constant $v$ still discards the temporal profile. The following
screen quantifies that limitation independently of any numerical example.

\begin{proposition}[A small included forcing value limits the constant barrier]
\label{cmp:small-slice}
Suppose the initial gap square satisfies $D(0)\le d_+$ and an included
slice covered by the proposed forcing floor has $K\le k_*<1$. Every
certificate using the constant parameter of Theorem~\ref{cmp:stop} then
has $v\ge\sqrt{1-k_*}$ and can satisfy its strict margin only if
\begin{equation}
 \Delta<\frac{d_+}{2\sqrt{1-k_*}}.
 \label{cmp:duration-cap}
\end{equation}
In particular, at a slice with $|D|\le\delta$ and
$|\mathcal F|\le A$, one may take $k_*=2\delta^2 A<1$.
\end{proposition}
\begin{proof}
A valid floor must have $\kappa\le K\le k_*$. The definition of $v$
therefore gives the asserted lower bound. Dropping the positive stop
threshold and enlarging the initial lower anchor to the actual initial
value, then to $d_+$, can only enlarge the allowable duration. This gives
\eqref{cmp:duration-cap}. The last statement is $|K|=2D^2|\mathcal F|$.
\end{proof}

This is a limitation of the displayed constant-$v$ consumer, not of all
comparisons using a constant $K$ bound, and not of actual dynamics.
Retaining the initial transient, using a time-varying barrier, renewing
anchors or integrating the regular $(D,M)$ system changes the consumer.
The supporting example in Appendix~\ref{app:compact-example} applies the
proposition to inherited certified data; it is not needed for the theorem.

An independent rectangular forcing floor can lose even the equilibrium
advantage. At a start with $V_0=0$, suppose the chosen reduction uses
$\mathcal F\ge f$ with $0\le f\le1/(2d_0^2)$ and
$d_{\mathrm{stop}}=\alpha d_0$, $0<\alpha<1$. It replaces the trial
forcing by $\kappa=2d_{\mathrm{stop}}^2f\le\alpha^2$, so
$v\ge\sqrt{1-\alpha^2}$. The exit margin requires
\begin{equation}
 \Delta<\frac{d_0}{2}\sqrt{\frac{1-\alpha}{1+\alpha}}<\frac{d_0}{2}.
 \label{cmp:rectangular-screen}
\end{equation}
Thus that independent reduction can be worse than the gap-only deadline.
It does not exclude correlated trial regions or a source estimate for
the combined kernel itself.

Every proposed source acquisition must also supply the direct comparator
with the same information: the centered discriminant, complete $M$
transport and regular $(D,M)$ evolution. If those shared coordinates
already prove positivity, computing $K$ adds a magnitude or cost
comparison, not new zero coverage. Likewise the lattice block identity
of Section~\ref{sec:lattice} does not by itself prevent concentration
into a single small gap. These are distinct limitations of scalar and
aggregate consumers, not reasons to discard the shared-source structure.
